\subsection{Produktivitätswirkung}
\label{subsec:produktivitaetswirkung}

Die Bewertung stellt Anforderungen, implementierte Mechanismen und Evidenz
gegenüber. Kriterien sind Wiederverwendung, Standardisierung, Variabilität,
Wartbarkeit, Testbarkeit, Reproduzierbarkeit, Entwicklerfreundlichkeit,
Erweiterbarkeit und Plattformabdeckung. Komplexität, Betriebsreife und
Dokumentierbarkeit ergänzen die Einordnung. Das projektinterne
Abnahmeprotokoll liefert Projektnachweise, aber keine unabhängige Replikation
\parencite{kleiveist2026acceptance}.
Damit bleiben beobachtete Mechanismen, ausgeführte Nachweise und abgeleitete
Wirkungen getrennt. Automatisierung gilt nicht allein deshalb als messbarer
Produktivitätsgewinn.

Produktivität umfasst hier den Aufwand bis zum lauffähigen Ausgangsprojekt,
wiederkehrende Einrichtungsentscheidungen, den Entwicklungsfluss und die
spätere Pflege. Eine einzelne Aktivitätskennzahl bildet sie nicht angemessen ab
\parencite{forsgren2021space}. Konzeptionell lässt sich die Wirkung daher als

\begin{center}
  \emph{Produktivitätsnutzen = vermiedener Projektaufwand
  $-$ Template-Pflegeaufwand $-$ zusätzliche Komplexität}
\end{center}

auffassen. Mangels Zeit- und Aufwandsmessungen ist sie nicht numerisch
auswertbar.
Die Beziehung zeigt zugleich, dass zentrale Pflege und zusätzliche Komplexität
den vermiedenen Projektaufwand mindern können.

Beim Projektstart bündelt \texttt{control.py} Profilwahl, Generierung,
Systemprüfung und Installation. Die Abnahme belegt diese Schritte für alle
fünf Profile sowie die Ablehnung inkompatibler PostgreSQL-Auswahlen
\parencite{kleiveist2026acceptance}. Die gemischten CI-Ergebnisse des späteren
HEAD begrenzen die Übertragbarkeit. Weniger manuelle Einrichtung ist daher
plausibel, eine konkrete Zeitersparnis aber nicht belegt.

In der Entwicklung vereinheitlichen profilabhängige Einstiege Start, Test und
Build. Gemeinsames Frontend, Konfigurationsvertrag und klare Verantwortungen
fördern Wiederverwendung; Fachwerkzeuge und Plattformvoraussetzungen bleiben.
Governance~v1 ergänzt einen einheitlichen, früheren Rückkanal für definierte
Größen-, Komplexitäts- und Abhängigkeitsbefunde. Zentrale Rule IDs und
Handlungshinweise können die Zuordnung eines Befunds für Menschen und Coding
Agents erleichtern. Gemessen wurden jedoch weder kürzere
Fehlerbehebungszeiten noch weniger Reviewaufwand. Die 42 Warnungen und acht
starken Warnungen des lokalen Reports sind ein Snapshot der Policy-Anwendung,
kein Produktivitätsmaß.

Zentrale Änderungen nützen künftigen Projekten, erfordern aber gemeinsame
Pflege von Profilen, Generator, Tests und Dokumentation. Bereits generierte
Projekte werden nicht automatisch aktualisiert.
Der Nutzen verteilt sich somit über Projektstart, laufende Entwicklung und
spätere Baseline-Pflege. Er ist besonders plausibel, wenn mehrere technisch
verwandte Vorhaben dieselben Grundlagen verwenden.

Quantitative Evidenz fehlt. Eine spätere kontrollierte Evaluation müsste
vergleichbare Projekte mit und ohne Template anhand von Start- und
Onboarding-Zeit, Setup-Fehlern, fehlgeschlagenen Builds, Wiederverwendungsanteil
und Pflegezeit untersuchen.
Zusätzlich wären die Zeit bis zur ersten Fachfunktion und die Zahl manueller
Einrichtungsentscheidungen zu erfassen. Erst solche Vergleichsdaten erlauben
eine Aussage über die Größe des Effekts.
\beginnerinput{workspace/statement/700_bewertung/710}
